\chapter{Problem Formulation}

Medical exam preparation is a laborious task requiring extensive rote practice to internalise clinical vocabulary, symptom presentations, conditions, risk thresholds, and pharmacological details. Independent study is difficult without feedback, peer interaction, or access to professional tutoring. These are resources that are unevenly distributed across the candidate population. \gls{emma} addresses this by providing an accessible, textbook-grounded conversational study agent capable of both answering clinical questions and administering adaptive quiz sessions. 

This project builds directly on the \gls{amg-rag} methodology introduced by \citet{rezaei2025}, using the same MedQA-\gls{usmle} and MedMCQA datasets. The choice to frame \gls{emma} as a study agent rather than a triaging system reflects the nature of these datasets: their multiple-choice structure, specialty labels, and associated textbook corpus are ideally suited to exam preparation rather than clinical decision support.

\section{Vocabulary and Term Definitions}

Before any \gls{ml} component is applied, the input space must be anchored to clinically valid terminology. Patient and student queries may use lay language (``chest hurts'', ``can't breathe'') that must be mapped to accepted medical terms (e.g., chest pain, dyspnea).

We define the accepted vocabulary $\Sigma$ as the set of clinical entity spans recognisable by the \texttt{en\_ner\_bc5cdr\_md} biomedical \gls{ner} model \cite{neumann2019}, specifically, spans labelled as \texttt{DISEASE} or \texttt{CHEMICAL} under the BC5CDR annotation schema. This model was trained on 1,500 PubMed articles and recognises 5,818 disease and 4,409 chemical entity types. \texttt{en\_core\_sci\_md} was replaced and not used for entity extraction since it outputs a single generic `ENTITY' label unsuitable for typed retrieval rewriting. 

Vocabulary coverage is extended by a \gls{pmi}-derived thesaurus mapping lay terms to clinical co-occurrents, constructed from collocation analysis of the MedQA training corpus (Chapter \ref{ch:ner_collocation}). Only tokens belonging to $\Sigma$ are extracted and used for \gls{faiss} query rewriting. Input text is otherwise passed to the \gls{llm} in full. This design separates the retrieval query (entity-focused, precise) from the generation prompt (full context, natural language).

\section{Research Question}

The central research question is: does textbook-grounded \gls{rag} improve small, locally-run \gls{llm} accuracy on MedQA compared to a no-retrieval baseline, and under what conditions? 

This question is addressed by a six-run ablation grid varying embedding model (Qwen3-Embedding-0.6B, Octen-Embedding-0.6B), \gls{llm} (Qwen3-4B, Qwen3-4B-Thinking-2507), and \gls{rag} condition (on/off), evaluated on 100 MedQA questions per run. Results are compared against the \gls{amg-rag} reference of 74.1\% F1 on MedQA using GPT-4o-mini and a dynamic PubMed knowledge graph \cite{rezaei2025}.

\section{Architecture}

\gls{emma} processes each user query through a five-stage pipeline before returning a response:

\begin{enumerate}
    \item \texttt{en\_ner\_bc5cdr\_md} extracts \texttt{DISEASE} and \texttt{CHEMICAL} entities from the raw query. These entities replace the full vignette text as the \gls{faiss} retrieval query, preventing incidental clinical language from diluting the embedding. If no entities are found (5.8\% of MedQA questions), the raw query is used unchanged. 

    \item The rewritten query is embedded and searched against a \gls{faiss} \gls{ifip} of 36,723 textbook chunks. Three embedding models were evaluated; Octen-Embedding-0.6B produced the best ablation results and is used in production. Retrieved chunks are assigned confidence bands across three tiers: high ($\geq$ 0.70), medium (0.55-0.70, flagged in the prompt with a hedging instruction), and low (0.40-0.55, included but marked uncertain). Chunks scoring below 0.40 are dropped entirely. 

    \item The raw query is independently passed to the \gls{tfidf} + \gls{lsvc} classifier, which assigns one of 19 specialty labels. This label is included in the \gls{llm} prompt as a routing context. The BERTopic clustering model provides finer-grained topic context within specialties. 

    \item All retrieved passages, entity spans, specialty labels, and confidence flags are assembled into a structured prompt with hedging instructions for low-confidence sources. 

    \item Inference is attempted via Ollama. This enables fast local paths, with no \gls{gpu} required. If Ollama is unavailable or the model is not pulled, the system falls back to HuggingFace transformers with \gls{nf4} quantization on a Colab T4 \gls{gpu}. The production \gls{llm} is Qwen3-4B-Thinking-2507, selected by ablation: the standard Qwen3-4B variant cannot effectively use retrieved context ($-4$\gls{pp} with \gls{rag}), whereas the Thinking variant delivers +11\gls{pp} with the Octen embedding. 
\end{enumerate}

The \gls{cf} recommender operates as a separate offline module, tracking per-specialty quiz accuracy and generating study recommendations independently of the \gls{rag} pipeline. See Figure \ref{fig:pipeline} in Appendix for Architecture Pipeline.

\section{Model Selection Rationale}
\label{sec:model_selection}

\subsection{Embedding Models}


Model selection followed the \href{https://huggingface.co/spaces/mteb/leaderboard}{MTEB Leaderboard}, specifically the RTEB Healthcare benchmark. 

\begin{enumerate}
    \item \href{https://huggingface.co/Qwen/Qwen3-Embedding-0.6B}{Qwen3-Embedding-0.6B} was selected since its 8B variant, Qwen3-Embedding-8B, ranks \#6 overall, the highest among models under 14B parameters (score 67.03). The 0.6B variant was chosen as the \gls{vram}-feasible member of that family given the Colab T4 budget. 
    
    \item \href{https://huggingface.co/Onten/Octen-Embedding-0.6B}{Octen-Embedding-0.6B} was selected because the Octen-8B variant ranks \#7 on the same leaderboard across all model sizes (score 73.79), outperforming Qwen3-Embedding-8B despite being in the same size class. Importantly, Octen-8B is derived from Qwen3-Embedding, meaning the 0.6B members of both families are directly comparable since the Octen variant is a domain-refined descendant. This made the two a natural ablation pair: same architecture family, different training regimen. 
\end{enumerate}

\subsection{LLM Models}
Model selection followed the \href{https://medmarks.ai}{Medmarks.ai} benchmark (last updated December 2025), filtered to models under $\sim$5B parameters.

\begin{enumerate}
    \item \href{https://huggingface.co/Qwen/Qwen3-4B-Thinking-2507}{Qwen3-4B-Thinking-2507} ranked \#33 with a mean win rate of 0.4918 and weighted mean of 0.4883, substantially above all other sub-5B models evaluated.
    \item The standard \href{https://huggingface.co/Qwen/Qwen3-4B}{Qwen3-4B} was included as a controlled baseline to isolate the effect of chain-of-thought reasoning, holding architecture and parameter count constant.
\end{enumerate}